\section{Situación actual}
\label{sec:intro}

Una comercializadora de energía ya usa IA en varias áreas: un asistente de facturación en servicio al cliente, mensajes
automáticos en cobranzas y resúmenes de llamadas en ventas. Cada uno funciona, pero cada uno
se armó en una herramienta distinta, con su propio modelo, sus propias llaves y su propio
acceso a datos de clientes (\cref{fig:islas}). El resultado es que nadie tiene la vista de
conjunto.

Esto tiene nombre en ingeniería de ML: cada agente aislado es un \emph{consumidor no
declarado} de datos \citep{sculley2015hidden}, algo que usa información sensible sin que la
organización lo sepa. Y el costo de no resolverlo es concreto: Gartner proyecta que más del
40\,\% de los proyectos agénticos se cancelarán antes de 2027 por costos que se salen de
control, valor poco claro o controles de riesgo insuficientes \citep{gartner2025agentic}. Las
tres causas son exactamente lo que una capa común vuelve visible y controlable.

\begin{figure}[!htbp]
  \centering
  \begin{tikzpicture}[x=1cm,y=1cm]
% ---------------- Izquierda: hoy ----------------
\node[titulo] at (3.4,4.55) {HOY: TRES ISLAS};
\foreach \i/\area/\ag/\herr/\mod in {0/Servicio al cliente/Asistente de\\facturación/Herramienta A/Modelo X,
                                     1/Cobranzas/Mensajes\\automáticos/Herramienta B/Modelo Y,
                                     2/Ventas/Resúmenes de\\llamadas/Herramienta C/Modelo Z}{
  \node[neut, minimum width=2.05cm, minimum height=0.95cm, font=\sffamily\scriptsize\bfseries] (a\i) at ({1.1+2.3*\i},3.4) {\ag};
  \node[plain, text=docgris] at ({1.1+2.3*\i},4.15) {\area};
  \node[risk, minimum width=2.05cm, font=\sffamily\tiny] (p\i) at ({1.1+2.3*\i},2.05) {\herr\\\mod\ · llaves propias};
  \draw[flecha] (a\i) -- (p\i);
}
\node[dat, minimum width=6.6cm, minimum height=0.6cm] (dc) at (3.4,0.85) {Datos de clientes (facturas, consumo, contacto)};
\foreach \i in {0,1,2}{ \draw[flechad, draw=riskbd] (p\i) -- node[etq, text=riskbd]{¿?} (p\i |- dc.north); }
\node[plain, text=riskbd, font=\sffamily\scriptsize\itshape] at (3.4,0.05)
  {¿Cuántos hay? ¿Qué modelo? ¿Qué datos? ¿Cuánto cuestan?};

% ---------------- Flecha central ----------------
\draw[-{Stealth[length=9pt,width=9pt]}, line width=2.2pt, draw=docacento] (7.15,2.3) -- (8.35,2.3);

% ---------------- Derecha: propuesta ----------------
\node[titulo] at (12.25,4.55) {PROPUESTA: UNA PLATAFORMA};
\foreach \i/\ag in {0/Asistente de\\facturación, 1/Mensajes de\\cobranza, 2/Resúmenes de\\llamadas}{
  \node[exec, minimum width=2.05cm, minimum height=0.95cm, font=\sffamily\scriptsize\bfseries] (b\i) at ({9.95+2.3*\i},3.4) {\ag};
  \node[plain, text=docgris, font=\sffamily\tiny] at ({9.95+2.3*\i},4.15) {manifiesto v1.x};
}
\node[ctrl, minimum width=6.6cm, minimum height=0.42cm] (r1) at (12.25,2.45) {\textbf{Registro} · inventarios de agentes, modelos y orquestación};
\node[ctrl, minimum width=6.6cm, minimum height=0.42cm, below=0.08cm of r1] (r2) {\textbf{Compuerta de evaluación} · nada se publica sin pasar};
\node[exec, minimum width=6.6cm, minimum height=0.42cm, below=0.08cm of r2] (r3) {\textbf{Gateway} · alias de modelo, costo, enmascaramiento};
\node[obs,  minimum width=6.6cm, minimum height=0.42cm, below=0.08cm of r3] (r4) {\textbf{Trazas y costos} por agente y por área};
\foreach \i in {0,1,2}{ \draw[flecha] (b\i) -- (b\i |- r1.north); }
\node[dat, minimum width=2.05cm, font=\sffamily\tiny] at (9.95,0.05) {Catálogo de\\modelos};
\node[dat, minimum width=2.05cm, font=\sffamily\tiny] at (12.25,0.05) {Herramientas\\MCP};
\node[dat, minimum width=2.05cm, font=\sffamily\tiny] at (14.55,0.05) {Conocimiento\\(RAG)};
\draw[flecha] (r4.south) -- (12.25,0.38);
\end{tikzpicture}

  \caption{De islas a plataforma: los mismos agentes, declarados en un registro y ejecutados a
  través de controles comunes.}
  \label{fig:islas}
\end{figure}

El documento se organiza en torno a seis preguntas de diseño. Cada sección abre con la respuesta corta y
cada decisión importante va acompañada de su razón y de la alternativa que descarté. Varias
de esas decisiones vienen de errores que ya vi en sistemas agénticos que construí y operé en
producción; cuando es así, lo digo.
